% Topic document -- 2 to 4 pages in the course template format.
\documentclass[11pt]{article}
\usepackage[utf8]{inputenc}
\usepackage[T1]{fontenc}
\usepackage[margin=1in]{geometry}
\usepackage{amsmath,amssymb}
\usepackage{booktabs}
\usepackage[round]{natbib}
\usepackage{xcolor}
\usepackage[colorlinks=true,allcolors=blue!50!black]{hyperref}

\newcommand{\kb}{k_B}
\newcommand{\km}{k_M}
\newcommand{\Rb}{R_B}
\newcommand{\Rm}{R_M}

\title{Project Topic: Credit Quantity, Credit Prices, and Future Output\\
\large A model of heterogeneous financing access}
\author{Alvaro Marcelo Ch\'avez Unyen\\
\small\url{https://github.com/amchavezu/ai-project}}
\date{Topic presentation, October 7, 2026}

\begin{document}
\maketitle

\section{The question and the track}

This is a \textbf{Track B} project: I formulate the economic model that is
currently missing from my undergraduate thesis. The question is: \emph{when
does the quantity of bank credit contain more information about firms' future
production than the price of credit, and can heterogeneous access to outside
finance explain why those margins may differ across sectors?}

The thesis is still in progress. Its preliminary forecasting estimates suggest
that credit quantities may be more informative for non-tradable activity at
relatively short horizons, while a credit-price indicator may become more
informative for tradable activity over a longer horizon. I treat this as a
qualitative pattern to motivate the model, not as a settled empirical fact or a
causal estimate. The project asks for a mechanism that can generate the
cross-sector ranking before trying to explain its exact timing.

\section{The baseline and the assumption I add}

The empirical thesis compares autoregressive forecasts with and without credit
indicators. That exercise can establish incremental predictive content, but it
does not contain a firm problem that distinguishes the quantity of funds from
their marginal price. I add firms that differ in their access to alternative
finance.

There are two dates. Firm $i$ enters date zero with internal funds $n_i>0$,
chooses capital $k_i$, bank borrowing $b_i$, and outside finance $m_i$, and
produces $F_i(k_i)$ at date one. It solves
\begin{equation}\label{eq:problem}
 \max_{k_i,b_i,m_i}\quad
 \beta F_i(k_i)-\Rb b_i-\Rm m_i
 \quad\text{s.t.}\quad
 k_i=n_i+b_i+m_i,\quad 0\leq b_i\leq\bar b_i,\quad m_i\geq0.
\end{equation}
For a bank-dependent firm, $m_i=0$. A firm with market access may use $m_i>0$
after exhausting its bank line. The maintained conditions are
\begin{align*}
 &F_i'(k)>0,\qquad F_i''(k)<0,\\
 &0<\Rb\leq\Rm=\Rb+\tau_i,\qquad \tau_i\geq0.
\end{align*}
The financing order means that the firm uses internal funds first, bank credit
second, and outside finance last. The bank line creates a quantity margin;
$\Rm$ creates a marginal-price margin. For closed forms and formal verification
I use
\begin{equation}\label{eq:technology}
 F_i(k)=a_i k-\frac{c_i}{2}k^2,\qquad a_i,c_i>0,
\end{equation}
on its increasing region. Let $\bar k_i=n_i+\bar b_i$ denote capital at the bank
ceiling. This is a deliberately small partial-equilibrium model. It takes the
credit line and finance wedges as given rather than claiming to identify the
source of a credit shock.

\section{The result I expect}

The left and right marginal conditions are
\begin{equation}\label{eq:focs}
 \beta F_i'(k)=\Rb
 \quad\text{and}\quad
 \beta F_i'(k)=\Rm.
\end{equation}
Under equation~\eqref{eq:technology}, the corresponding targets are
\[
 \kb=\frac{\beta a_i-\Rb}{\beta c_i},
 \qquad
 \km=\frac{\beta a_i-\Rm}{\beta c_i},
 \qquad \km\leq\kb.
\]
I expect the market-access firm's policy to be
\begin{equation}\label{eq:policy}
k_i^*=\begin{cases}
n_i, & \kb\leq n_i,\\
\kb, & n_i<\kb<\bar k_i,\\
\bar k_i, & \km\leq\bar k_i\leq\kb,\\
\km, & \bar k_i<\km.
\end{cases}
\end{equation}

The main conjecture compares the last two economically relevant cases, holding
the regime fixed under a small parameter change. If
$\km<\bar k_i<\kb$, the bank ceiling binds at a kink and
\[
 \frac{\partial k_i^*}{\partial\bar b_i}=1,
 \qquad
 \frac{\partial F_i(k_i^*)}{\partial\bar b_i}=F_i'(\bar k_i)>0.
\]
The available quantity of bank finance therefore determines the scale of the
project. A small common spread movement has no local real effect while the firm
stays at the kink.

If $\bar k_i<\km$, outside finance funds the marginal unit. A larger bank line
then replaces outside finance without changing total capital, while a common
credit-cost factor $z$ that raises both $\Rb$ and $\Rm$ gives
\[
 \frac{\partial k_i^*}{\partial\bar b_i}=0,
 \qquad
 \frac{\partial k_i^*}{\partial z}
 =\frac{1}{\beta F_i''(\km)}<0.
\]
The spread matters because it changes the cost of the marginal unit, not merely
because it correlates with the business cycle.

Let $\lambda_s$ be the share of firms in sector $s$ at the quantity-sensitive
kink. If $\lambda_N>\lambda_T$ and marginal products are positive, then the
non-tradable sector is more sensitive to the credit ceiling and the tradable
sector is more sensitive, in absolute value, to the common spread. This is the
qualitative bridge to the preliminary thesis pattern. Slow-moving financing
access can make current credit quantities reveal the near-term investment
pipeline of constrained firms. Credit prices can contain information about the
desired scale of projects for firms that can substitute across funding sources.
Different adjustment speeds may move the peak response across horizons, but the
baseline model does not yet derive exact forecast horizons.

\section{Why this is not already solved}

I compared the published articles, publicly available working-paper versions,
and appendices of the closest contributions. \citet{holmstromtirole1997}
derive how firm and intermediary capital determine monitored and direct
finance. Their mechanism motivates heterogeneous access, but it does not rank
credit quantities and prices as sector-specific forecasting indicators.
\citet{gilchristzakrajsek2012} show that corporate-bond spreads predict
aggregate activity and decompose the spread into expected-default and excess
premium components. Their working-paper forecasting exercises cover short and
longer horizons, but they do not compare that price signal with a binding bank
quantity at the firm level.

\citet{mullerverner2024} are the closest sectoral paper. Their published
\emph{Review of Economic Studies} article and its NBER working-paper appendix
show that non-tradable credit expansions have different macroeconomic content
from tradable credit expansions and examine bond issuance as a robustness
check. Their object is the allocation of credit across sectors, not the choice
between a capped bank tranche and a more expensive marginal source. Follow-on
work such as \citet{gilchristsimzakrajsek2014} models how spreads affect firm
investment, but does not deliver the quantity-versus-price sectoral ranking.
The proposed contribution is therefore narrow: a transparent threshold model
tailored to the thesis's two predictors. I do not claim that credit rationing,
outside finance, or sectoral heterogeneity is new.

\section{Plan and risks}

The next steps are to (i) prove the complete piecewise policy and both sectoral
comparative statics; (ii) make \texttt{code/verify.py} reproduce every reported
number and fail outside the stated conditions; (iii) test behavior close to the
regime boundaries; and (iv) formalize every numbered result in Lean using the
course's AppliedModelingLib workflow. A dynamic extension will add capital
accumulation, adjustment costs, and the shadow value of the credit limit, but it
will remain separate from the baseline proposition until its conditions can be
proved.

The main empirical risk is that the thesis data do not directly measure firms'
outside-finance access, so the sectoral composition condition
$\lambda_N>\lambda_T$ may remain a plausible hypothesis rather than an
identified fact. The main theoretical risk is that the sharp local zero in the
kink regime disappears once the credit ceiling or funding costs are jointly
endogenous. If either risk materializes, I will keep the result conditional,
report regime switching in simulations, and avoid using the model to claim
exact timing or causality. The Lean risk is the piecewise optimum; the fallback
is to formalize the quadratic specialization first and label any remaining
general concave result as open.

\bibliographystyle{plainnat}
\bibliography{../paper/references}

\end{document}
